% Pursuing Stacks (cotes 134-2 to 134-8): archive cote 134-6, site page 16 (PDF page 17), typescript page 335.
% typescript. Draft compared: ps5.tex:2169-2241.
% AI-transcribed by Claude Opus 5.5 (claude-opus-5-5, adaptive thinking, effort high), Message Batches API, batch ps3 (msgbatch_01XdFmuVVBAcgc7EFxdtCqsj), 2026-11-12; reviewed against both supplied scan halves; see CHECK_DECISIONS.tsv.
\scanpage{134-6/16}
\mspage{335}

signatures (13) by the ``Stokes rule'' (17), the corresponding operators
(11) do turn the family $(L_n)$ into a chain complex, na\over[typed]{\unsure{e}}{m}\over[typed]{\unsure{m}}{e}ly we have
the relations
\[ (18) \qquad d_{n-1}d_n = 0 \quad \hand{\unsure{,|.}} \]
\over{\unsure{It}}{This} follows immediately from the well-\over[typed]{\illegible{1}}{k}nown anticommutativity property
of \xout{induction} (twofold) induction of orientation on boundaries.

Things are a little more delicate if we do'nt assume that isomor-%
phisms are identities, even if (by compens\unsure{q|a}tion) we \inserted[typed]{should} insist that two
distinct objects are never isomorphic. To define $d_{\hand{n}}$ , we then must
\uline{cho\inserted[typed]{o}se}, for any subobject \over{b}{$\beta$} of dimension n-1 of \xout{the} an object \over{a}{$\unsure{a|\alpha}$} of
dimension n, \inserted[typed]{just \uline{one}} \xout{a} representative monomorphism (13) of \hand{$\beta$}. The coherence
condition then needed in order to get (18), is that \xout{for} any \xout{two such}
square diagram
% Hand-drawn arrows are marked blue; typed objects remain plain.
\[ (19) \qquad \diagram{\begin{tikzcd}[sep=small]
 & \over{\unsure{e|a}}{c} \ar[color=blue,dl] \ar[color=blue,dr] & \\
 b \ar[color=blue,dr] & & b\hand{'} \ar[color=blue,dl] \\
 & a &
\end{tikzcd}} \]
made up with such restricted boundary maps, should commute -- a some-%
what delicate condition, presumably hard to insure, for the choices
involved for defining the ``strict'' boundary maps in A.

In one case as in the other, we are very close of course of the
set up envisioned in section 91 - it would'nt be hard even to \xout{subordi\unsure{n}}
fit the case considered here into this set up, if we make the slight
extra assumption th\over[typed]{e}{a}\over[typed]{n}{t} any map in A factors into an ep\xout{o}imorphism-with-%
section, followed by a monomorphism (which is true indeed in the
``standard'' cases), which will ensure that for varying a, i(a) is indeed
a functor with values in (Ord), as stated in loc. cit. But from the
point of view of construction of abelianizators, it would seem that
the e\over[typed]{i}{x}\over[typed]{x}{i}\over[typed]{t}{s}\over[typed]{s}{t}\over[typed]{n}{e}\over[typed]{e}{n}ce of \over{a}{the} functor
\[ \hand{(20)} \qquad i: A \hand{\longrightarrow} (\mathrm{Ord}) \]
is irrelevant.

O\over[typed]{r}{u}r main question now, of course, is about \uline{the chain complex L.}
\uline{being an abelianizator or not}. The question is interesting even in the
standard cases, by cho\inserted{o}sing the orientations (16) in a way different
from the standard one. Are the corresponding chain\xout{s} complexes in
Add($A^{\mathrm{o}}$) necessarily chain-homotopic ?

It just occurs to me that \inserted{indeed,} between the chain complexes

\margin{16}

% Visual collation: both supplied image halves inspected; see CHECK_DECISIONS_134_7_8.tsv.
